\section{多步梯度更新的问题}

On-policy actor-critic 每收集一批数据只能做一步梯度更新，效率很低。我们希望在同一批数据上做\textbf{多步}更新。

\subsection{Importance Weights 回顾}

用旧策略 $\pi_\theta$ 的数据来更新新策略 $\pi_{\theta'}$ 时，需要 importance weights：
$$\nabla_{\theta'} J(\theta') \approx \sum_{t,i} \frac{\pi_{\theta'}(a_{i,t} \mid s_{i,t})}{\pi_\theta(a_{i,t} \mid s_{i,t})} \nabla_{\theta'} \log \pi_{\theta'}(a_{t,i} \mid s_{t,i}) \hat{A}^{\pi_\theta}(s_{t,i}, a_{t,i})$$

\begin{warningbox}{多步更新的风险}
如果对代理目标函数（surrogate objective）做很多步梯度上升，策略会被激励大幅偏离旧策略。此时 importance weights 不再可靠，优势估计也失效，导致过拟合和性能崩塌。
\end{warningbox}

\subsection{本章小结}

多步梯度更新能提高数据效率，但需要约束策略更新的幅度来保持稳定性。

